\section*{Appendix A: Continuity Receipt Schema}
\label{app:schema}
\addcontentsline{toc}{section}{Appendix A: Continuity Receipt Schema}

The Continuity Receipt is the core protocol object of CPS-0001. Each receipt captures a bounded interval of observation and carries evidence from one or more engines.

\begin{verbatim}
ContinuityReceipt {
  protocolVersion: String,        // "1.0"
  receiptId: UUID,
  interval: { start: DateTime, end: DateTime, coverageMs: UInt64 },
  subject: { id: SHA256, type: "embodied" | "agent" },
  evidence: [{
    engineId: String,
    engineVersion: String,
    confidence: Float,            // 0.0 — 1.0
    payloadDigest: SHA256,
    payload: JSON
  }],
  assertions: {
    observationOccurred: { value: Bool, confidence: Float },
    continuityMaintained: { value: Bool, confidence: Float },
    receiptIntegrity: { value: true, confidence: 1.0 }
  },
  issuer: { id: String, publicKey: HexString },
  previousReceiptHash: SHA256?,   // null for chain head
  expiresAt: DateTime,
  signature: { algorithm: "Ed25519", value: HexString, signedAt: DateTime }
}
\end{verbatim}

\subsection*{Verification Contract (V$_1$–V$_7$)}

Each receipt is verified against seven checks:
\begin{enumerate}
\item[V$_1$] Signature validity (Ed25519)
\item[V$_2$] Receipt integrity (hashes match payloads)
\item[V$_3$] Temporal validity (not expired)
\item[V$_4$] Chain continuity (predecessor hash matches)
\item[V$_5$] Minimum confidence threshold
\item[V$_6$] Engine eligibility (known engine list)
\item[V$_7$] Subject consistency across chain
\end{enumerate}

A receipt passes iff all seven checks pass.

\newpage
\section*{Appendix B: Experimental Protocol}
\label{app:protocol}
\addcontentsline{toc}{section}{Appendix B: Experimental Protocol}

\subsection*{B.1 PES Benchmark (EE-001)}

\textbf{Hardware:} Standard webcam, 30 fps. MediaPipe Pose 33-landmark model.
\textbf{Environment:} Indoor, varied lighting (daylight, fluorescent, LED).
\textbf{Subjects:} 54 participants (self-selected, age 22–61).
\textbf{Duration:} 30 seconds per capture.
\textbf{Processing:} On-device browser-side only. No raw motion data transmitted.

\subsection*{B.2 Causal Coupling (EE-002)}

\textbf{Devices:} Phone IMU + webcam, time-synchronized via shared start signal.
\textbf{Metrics:} Temporal alignment precision at $\pm$16ms, 33ms, 66ms thresholds.
\textbf{Procedure:} Subject performs a sharp motion event (clap, nod, drop). Both sensors record independently. Post-hoc alignment computed via cross-correlation of acceleration and optical flow.

\subsection*{B.3 Challenge-Response (EE-003)}

\textbf{Prompt set:} 5 rotation directions (left, right, forward, back, shake).
\textbf{Timing:} Randomized jitter $\pm$200ms around 3s interval.
\textbf{Pass criterion:} Detected motion direction matches prompt direction within 1.5s.

\subsection*{B.4 Verification Session (VS-001)}

\textbf{Structure:} Passive presence detection (30s) → if pass, challenge-response (3 rounds).
\textbf{Multiple failure handling:} Any failed round falls back to re-authentication.
\textbf{N:} 60 sessions across 20 subjects, 3 sessions each.

\newpage
\section*{Appendix C: Threat Model}
\label{app:threat}
\addcontentsline{toc}{section}{Appendix C: Threat Model}

\begin{tabular}{|l|l|l|l|}
\hline
\textbf{Attack} & \textbf{Target} & \textbf{Feasibility} & \textbf{Mitigation} \\
\hline
Synthetic video frame injection & EE-001 & Medium (offline generation) & Cross-modal binding (EE-002) \\
Recorded motion replay & EE-003 & Low (requires jitter prediction) & Randomized timing \\
Sensor desynchronization & EE-002 & Medium (jam one sensor) & Chain continuity (V$_4$) \\
Credential replay & CPS-0001 & Low (signature binding) & V$_1$ + expiry (V$_3$) \\
Chain truncation & CPS-0001 & Medium (drop receipts) & V$_4$ predecessor check \\
Long-interval evasion & CPS-0001 & High (wait for expiry) & Requires active re-verification \\
\hline
\end{tabular}

\noindent\textbf{Key observation:} attacks that succeed against a single engine fail against the combined protocol. The forger must simultaneously defeat temporal, cross-modal, and challenge-response checks — each requiring a different attack capability.

\newpage
\section*{Appendix D: Published Failure Cases}
\label{app:failures}
\addcontentsline{toc}{section}{Appendix D: Published Failure Cases}

Following the research principle that negative results are results, we document known failures:

\subsection*{D.1 FD-001: Frame Rate Hypothesis}

\textbf{Claim:} Higher capture frame rate (60 fps) should improve PES separation.
\textbf{Result:} Falsified. 60 fps introduced sensor artifacts that degraded all four entropy dimensions. The optimal rate for entropy extraction is 30 fps.
\textbf{Reference:} Research Note FD-001.

\subsection*{D.2 DL-001: Direction Asymmetry in EE-003}

\textbf{Observation:} Challenge-response pass rates are asymmetric by direction. Pitch (forward/back) passes more reliably than yaw (left/right).
\textbf{Hypothesis:} Gyroscope axis sensitivity differences. Not yet confirmed.
\textbf{Reference:} Research Note DL-001.

\subsection*{D.3 Environmental Artifact: Fluorescent Flicker}

\textbf{Observation:} Frequency entropy dimension (EE-001) produced false positives under fluorescent lighting. The 60 Hz AC flicker was initially interpreted as biological frequency structure.
\textbf{Fix:} Notch filter at 60 Hz before frequency-domain analysis.

\subsection*{D.4 Spline Interpolation Proximity (Ongoing)}

\textbf{Observation:} AI strategy using cubic spline interpolation achieved PES 0.38 vs. human floor 0.41. This 0.03 margin is the narrowest separation observed.
\textbf{Status:} Active research. Potential mitigations: cross-modal binding forces spatial consistency that splines cannot maintain.
